% from ega0-20-fr.tex
\oldpage[0\textsubscript{IV-1}]{138}
où l'opérateur différentiel est $v^0$ (avec les notations de (20.2.1)).
La proposition résulte alors de la suite exacte (20.2.4.2) et de la
définition des modules de cohomologie
\begin{equation}
\tag{20.6.6.3}
\label{0.20.6.6.3-fr}
  \mathrm{H}^i(\mathrm{K}_\bullet,L)=
  \mathrm{H}^i(\operatorname{Hom}_C(\mathrm{K}_\bullet,L)).
\end{equation}
\end{proof}

Si on a un diagramme commutatif d'homomorphismes d'anneaux
\[
\xymatrix{
  A' \ar[r] & B' \ar[r] & C' \\
  A \ar[u] \ar[r] & B \ar[u] \ar[r] & C \ar[u]
}
\]
les di-homomorphismes (20.5.4.3) définissent un di-homomorphisme de
complexes de modules
\begin{equation}
\tag{20.6.6.4}
\label{0.20.6.6.4-fr}
  \mathrm{K}_\bullet(C/B/A)\longrightarrow
  \mathrm{K}_\bullet(C'/B'/A')
\end{equation}
et les di-homomorphismes qu'on en déduit pour l'homologie ou la
cohomologie s'identifient, par les formules (20.6.5.2), (20.6.6.1) et
(20.6.6.2), aux di-homomorphismes déjà définis dans (20.5.4.3),
(20.6.4.2), (20.2.1) et (18.3.7.2) (compte tenu, pour le dernier, de la
définition de l'opérateur $\partial$ dans la suite exacte (20.2.4.2)).

\begin{env}[20.6.7]
\label{0.20.6.7-fr}
On sait que pour un complexe $\mathrm{K}_\bullet$ de $C$-modules et un
$C$-module $L$, on a des homomorphismes canoniques
\[
  \alpha_i:\mathrm{H}^i(\mathrm{K}_\bullet,L)\longrightarrow
  \operatorname{Hom}_C(\mathrm{H}_i(\mathrm{K}_\bullet),L)
\]
(M, IV, 6). Ici, l'homomorphisme canonique
\[
  \alpha_1:\mathrm{H}^1(\mathrm{K}_\bullet(C/B/A),L)
  \longrightarrow
  \operatorname{Hom}_C(\mathrm{H}_1(\mathrm{K}_\bullet(C/B/A)),L)
\]
se définit immédiatement comme obtenu par passage au quotient par l'image
de $\operatorname{Hom}_C(\mathrm{K}_0,L)$ de l'homomorphisme de restriction
\[
  \operatorname{Hom}_C(\mathrm{K}_1(C/B/A),L)\longrightarrow
  \operatorname{Hom}_C(\mathrm{H}_1(\mathrm{K}_\bullet(C/B/A)),L)
\]
puisque $\mathrm{H}_1(\mathrm{K}_\bullet(C/B/A))$ n'est autre que le
noyau de $\mathrm{K}_1\to\mathrm{K}_0$; tenant compte de (20.6.6.2) et
(20.6.5.2), on obtient donc un $C$-homomorphisme canonique
\begin{equation}
\tag{20.6.7.1}
\label{0.20.6.7.1-fr}
  \operatorname{Exalcom}_{B/A}(C,L)\longrightarrow
  \operatorname{Hom}_C(\Upsilon_{C/B/A},L)
\end{equation}
qui s'explicite de la façon suivante : en vertu de (20.2.4.2), toute
$B$-extension de $C$ par $L$ qui est $A$-triviale provient de la donnée
d'une $A$-dérivation $D$ de $B$ dans $L$, donc (20.4.8) d'un
$C$-homomorphisme $f$ de $\Omega^1_{B/A}\otimes_B C$ dans $L$; on fait
correspondre à la classe de cette extension la \emph{restriction} de $f$
à $\Upsilon_{C/B/A}$, qui ne dépend effectivement que de cette classe et
non du choix de $D$.
\end{env}

\begin{definition}[20.6.8]
\label{0.20.6.8-fr}
Soient $u:A\to B$, $v:B\to C$ deux homomorphismes d'anneaux. Pour toute
$B$-extension $E$ de $C$ par un $C$-module $L$, qui est $A$-triviale
(18.3.7), on appelle \emph{homomorphisme caractéristique de $E$}, et on
note $\chi_E$, le $C$-homomorphisme $\Upsilon_{C/B/A}\to L$, image de la
classe de $E$ par l'homomorphisme canonique (20.6.7.1).
\end{definition}

